(1) 单重态的约化密度矩阵是 \(\rho_1=\rho_2=I/2\)，所以任意局域 Pauli 方向的平均值都为零。利用 \(\langle\psi^-|\sigma_{1i}\sigma_{2j}|\psi^-\rangle=-\delta_{ij}\)，得到
\[
E(a,b)=\langle(a\cdot\sigma_1)(b\cdot\sigma_2)\rangle=-a\cdot b.
\]
若两方向相同，测量结果总是反相关。

(2) 测量结果 s、t 的投影算符分别为 \([I+s(a\cdot\sigma)]/2\) 及 \([I+t(b\cdot\sigma)]/2\)，故
\[
P(s,t|a,b)=\frac14[1-st(a\cdot b)].
\]
四个结果的概率非负且和为一；对任一端结果求边缘概率始终为1/2。

(3) 定义 \(S=E(a,b)+E(a,b')+E(a',b)-E(a',b')\)。选择
\(a=(1,0),a'=(0,1),b=(1,1)/\sqrt2,b'=(1,-1)/\sqrt2\)，则三个加项均为 \(-1/\sqrt2\)，减掉的第四项为 \(+1/\sqrt2\)，因此
\[
S=-2\sqrt2,\qquad |S|=2\sqrt2.
\]
该值超出经典局域界2，但未超过量子 Tsirelson 上限。

(4) 对经典局域隐变量，每次试验中预先确定的四个输出 \(A,A',B,B'\) 均为 ±1，那么 \(A(B+B')+A'(B-B')=\pm2\)；取平均得到 |S|≤2。量子单重态的某些角度组合不满足此界，但每一端单独观察的结果概率仍为1/2，不依赖远端选择了 b 还是 b'，因而没有可控的超光速信号。
